\section{Nullability for a Graph You Share}
\label{sec:ch14-nullability-for-a-graph-you-share}

\index{nullability!the federated rule|(}%
Chapter~\ref{ch:schema-design} gave a rule: mark a field non-null when you can
guarantee it from data you own, in the same process, with a constraint behind
it. Every field in this chapter that emptied something passed that rule.
\code{ratingCount} is a count over a table its service owns. It is still the
field that emptied the response.

\index{nullability!the second question}%
So the rule is necessary and it is not sufficient, and the half it is missing is
not about your service at all. The rule I use now has a second clause on it:
mark a field non-null when the value is guaranteed from data you own, in the
same process, with a constraint behind it, \emph{and} the position it occupies
in the composed schema is one you are willing to empty when your service is
down.

\index{nullability!the advice that predates federation}%
Half of that is old advice. Sashko Stubailo argued in 2018, before any of this
tooling existed, that a field whose resolver fetches data asynchronously should
be nullable, so that a service or database being unreachable is easier to deal
with~\autocite{stubailo2018}. He was right, and he was writing about one
server. The second clause is mine, and I state it as mine because I went
looking for a published rule about the federated case and did not find one that
clears this book's bar for a source.

\index{nullability!two positions decide everything}%
The positions that matter are the two this chapter has met. A field contributed
to a type another service owns sits inside that service's list, and emptying it
empties the list. A root field sits directly under the response, and emptying it
empties the response. Everything else in a subgraph is comparatively cheap to
get wrong.

\subsection{What Moved, and What Did Not}

\index{nullability!the fields that changed}%
Five fields in this graph fail that second test, and all five give up their
exclamation marks. The root field goes first, because it is the one that costs
a whole response. Here is the Ratings service's \code{Query.cs} again,
complete, with what moved marked:

\begin{minted}[highlightlines={11,12,13,14,15,16,17,18,19}]{csharp}
using Microsoft.EntityFrameworkCore;

namespace Ratings;

[QueryType]
public static class Query
{
    /// <summary>
    /// How many scores this service holds, across every session.
    /// </summary>
    // The question mark is the whole of chapter 14. A count over this
    // service's own table is a value it can always produce, which is
    // what the nullability rule of chapter 4 asks of a non-null field.
    // Behind a router it is also a root field, and the nearest
    // position above a root field that is allowed to hold a null is
    // the response's own data entry. Declared non-null, this service
    // being stopped discards every other service's answer in the same
    // response along with it.
    public static async Task<int?> GetRatingCountAsync(
        RatingContext db,
        CancellationToken ct)
        => await db.Ratings.CountAsync(ct);
}
\end{minted}

\index{nullability!the two fields on somebody else's type}%
The same service contributes two more to a \code{Session} it does not own, and
they are a different case. This file is long enough that what follows is the
complete member rather than the whole of it, which is the other half of the
rule this book prints changes under:

\begin{minted}[highlightlines={4,5,6,7,8,9,10}]{csharp}
    /// <summary>
    /// How many scores this session has been left.
    /// </summary>
    // Nullable for a different reason than the root field above. This
    // one is contributed to a type another service owns, so it rides
    // inside that service's list. Non-null here does not empty the
    // response; it empties the page of sessions the Sessions service
    // answered correctly, because a list of non-null items has nowhere
    // to put one null and goes null entirely.
    public static async Task<int?> GetRatingCountAsync(
        [Parent(nameof(Session.Key))] Session session,
        IRatingsBySessionIdDataLoader ratingsBySessionId,
        CancellationToken ct)
        => (await ratingsBySessionId.LoadAsync(session.Key, ct))
            ?.Length ?? 0;
\end{minted}

\index{@requires@\code{@requires}!the field with two ways to fail}%
And the one beside it, which is the same change and the worse position, because
the router cannot ask for this field at all until it has collected a title from
a third service:

\begin{minted}[highlightlines={8,9,10,11,12,13,15}]{csharp}
    /// <summary>
    /// Where an attendee goes to leave a score for this session.
    /// </summary>
    // The url is built from the session's title, and the title belongs
    // to another service. @requires is what makes the router send it
    // along with the key, and [External] on GetTitle above is what
    // makes that declaration legal.
    //
    // Nullable on the same reasoning as ratingCount above, and the
    // @requires makes it the worse case rather than a better one: this
    // field cannot be answered unless the router has already collected
    // a title from a third service, so it has two ways to fail and
    // both of them land on somebody else's list.
    [Requires("title")]
    public static string? GetFeedbackUrl([Parent] Session session)
        => "https://feedback.example/2026/" + Slug(session.Title);
\end{minted}

\index{nullability!a method that cannot return null and a schema that says it can}%
The Search service's is the strangest of the five, because the method it sits
on cannot return null and the schema now says it can:

\begin{minted}[highlightlines={1,2,3,4,5,6,7,8}]{csharp}
    // This method cannot return null and the schema says it can, which
    // looks like an error and is the point. What the router puts in
    // this position is not what this method returns: it is whatever
    // the Sessions service answers for the key this method formats. A
    // non-null here is a promise about that service's uptime, made by
    // a service that has no way to keep it, and it costs the whole
    // search page when Sessions is stopped.
    public static Session? GetSession(
        [Parent] SessionSearchDocument document,
        INodeIdSerializer serializer)
        => new(serializer.Format(nameof(Session), document.SessionId));
\end{minted}

\index{nullability!what a resolver returns is not what the position holds}%
In one process the return type of a resolver is the type of the position it
fills, so a method that cannot return null describes a field that cannot be
null. Across a seam those come apart, and this method is where they do: it
hands the router a key, and what lands in the response is whatever another
service makes of that key.

\index{nullability!the fields that stay non-null}%
What does not move matters as much. A session's title, a speaker's name, and
the rating count and start time on a search document all keep their exclamation
marks. Each one is answered by the service that owns the row, out of a column
with a constraint behind it, in a position that same service is responsible
for. This is not an argument for a schema with no exclamation marks. It is an
argument for putting them where you can keep them.

\subsection{The Same Four Requests}

\index{partial failure!the graph after the change}%
Ratings stopped, and the request that lost the schedule:

\begin{minted}[fontsize=\footnotesize,breakanywhere]{json}
{"errors":[{"message":"Failed to fetch from Subgraph 'ratings' at Path 'sessions.nodes'."}],
"data":{"sessions":{"nodes":[
{"title":"Schemas That Outlive Their Authors","ratingCount":null},
{"title":"Reading a Query Plan","ratingCount":null},
{"title":"Paging Without Offsets","ratingCount":null},
{"title":"The Bank That Split Its Graph","ratingCount":null}]}}}
\end{minted}

\index{errors!one, because there is nothing to propagate}%
One error rather than two, everywhere in this section. The second error was
never a second failure; it was the propagation, and there is nothing left to
propagate.

\index{partial failure!the response survives}%
The root field, and the request that emptied everything:

\begin{minted}[fontsize=\footnotesize,breakanywhere]{json}
{"errors":[{"message":"Failed to fetch from Subgraph 'ratings'."}],
"data":{"ratingCount":null,"sessions":{"nodes":[
{"title":"Schemas That Outlive Their Authors"},
{"title":"Reading a Query Plan"},
{"title":"Paging Without Offsets"},
{"title":"The Bank That Split Its Graph"}]}}}
\end{minted}

\index{statement count!the same work, now delivered}%
The Sessions service spends the same one statement it always spent. The
difference is that the reader gets what it read.

\index{partial failure!the search page degrades}%
And the other direction across the same seam. Stop Sessions rather than
Ratings, and ask Search for the rating-descending page
chapter~\ref{ch:cross-seam-paging} built:

\begin{minted}[fontsize=\footnotesize,breakanywhere]{json}
{"errors":[{"message":"Failed to fetch from Subgraph 'sessions' at Path 'searchSessions.nodes.@.session'."}],
"data":{"searchSessions":{"nodes":[
{"session":null,"averageScore":4.5},
{"session":null,"averageScore":4}]}}}
\end{minted}

\index{partial failure!the ranking outlives the titles}%
The ranking survives without the titles. Search still spends its one statement,
the filter and the sort and the cut all happened, and what a client cannot do
is name the two sessions it found. That is a degraded search page rather than
no search page, and before this section it was no search page.

\subsection{Where This Leaves the Client}

\index{nullability!the bill the client pays}%
None of this is free, and the bill goes to the client.
Chapter~\ref{ch:schema-design} said so: make everything nullable and every field
of every client grows a null check, paid forever against failures that mostly do
not happen. This chapter has just made five fields nullable, and every client
that reads them now needs a branch it did not need before.

\index{nullability!the trade, stated plainly}%
That is the trade, and I take it. A null check is a line of client code. A
non-null field behind a seam is an outage in somebody else's service arriving
as an empty response in yours, at a moment nobody chose, with a status code of
200 on it. The first is annoying and the second is a page that does not render.

\index{nullability!what would remove the trade}%
What would be better is a schema that tells a client both things at once: this
value is always present, and it may be missing if something went wrong. That is
not a thing GraphQL can say today, and the next section is about the people
trying to make it one.
\index{nullability!the federated rule|)}%
